\documentclass[a4paper,11pt]{article}

% ─── Geometry ────────────────────────────────────────────────────────────────
\usepackage{geometry}
\geometry{margin=1in, headheight=15pt}

% ─── Mathematics ─────────────────────────────────────────────────────────────
\usepackage{amsmath}
\usepackage{amssymb}
\usepackage{mathtools}

% ─── Graphics / Colours ──────────────────────────────────────────────────────
\usepackage{graphicx}
\usepackage{xcolor}
\definecolor{blue3}{RGB}{0,0,180}
\definecolor{green3}{RGB}{0,120,0}
\definecolor{orange3}{RGB}{200,100,0}
\definecolor{red3}{RGB}{180,0,0}
\definecolor{grey3}{RGB}{80,80,80}

% ─── Units ───────────────────────────────────────────────────────────────────
\usepackage{siunitx}
\sisetup{detect-all}

% ─── Links ───────────────────────────────────────────────────────────────────
\usepackage[colorlinks=true,linkcolor=blue3,citecolor=blue3,urlcolor=blue3]{hyperref}
\usepackage{cleveref}

% ─── Boxes ───────────────────────────────────────────────────────────────────
\usepackage{tcolorbox}
\tcbuselibrary{skins,breakable}

\newtcolorbox{answerbox}[1][]{
    colback=red3!5,
    colframe=red3!60,
    coltitle=red3,
    title={Solution:},
    fonttitle=\bfseries,
    breakable,
    #1
}

% ─── Tables / Lists ──────────────────────────────────────────────────────────
\usepackage{booktabs}
\usepackage{enumitem}

% ─── Notation commands ───────────────────────────────────────────────────────
\newcommand{\ktilde}{\tilde{k}}
\newcommand{\ytilde}{\tilde{y}}
\newcommand{\ctilde}{\tilde{c}}
\newcommand{\dd}{\mathrm{d}}
\newcommand{\bgp}{\mathrm{BGP}}

% ─── Header / Footer ─────────────────────────────────────────────────────────
\usepackage{fancyhdr}
\pagestyle{fancy}
\fancyhf{}
\lhead{\textbf{14E022 Macroeconomics I} \quad Practice Exam 5 -- Solutions}
\rhead{BSE 2025--26}
\cfoot{\thepage}

% ═════════════════════════════════════════════════════════════════════════════
\begin{document}
% ═════════════════════════════════════════════════════════════════════════════

\begin{center}
    {\LARGE \textbf{14E022 Macroeconomics I}}\\[6pt]
    {\Large \textbf{Practice Exam 5 --- Answer Key}}\\[8pt]
    {\large BSE --- Academic Year 2025--26}
\end{center}

\noindent\rule{\textwidth}{0.4pt}

% ═════════════════════════════════════════════════════════════════════════════
\section{Question 1 --- Ramsey Model with Labour Income Tax}
% ═════════════════════════════════════════════════════════════════════════════

\subsection*{1.\ (6 points)}

\begin{answerbox}
\textbf{Step 1: Household budget constraint.}

The household receives after-tax labour income $(1-\tau)w_t$, capital income $r_t a_t$
(untaxed), and the lump-sum transfer $T_t = \tau w_t$. The flow budget constraint is:
\[
    \dot{a}_t = r_t a_t + (1-\tau)w_t + \tau w_t - c_t
    = r_t a_t + w_t - c_t.
\]
The tax and transfer cancel \emph{exactly}: the household's effective budget constraint is
identical to the no-tax case.

\textbf{Step 2: Current-value Hamiltonian.}

Let $a_t$ be the state variable and $c_t$ the control. The current-value Hamiltonian is:
\[
    \mathcal{H} = \frac{c^{1-\sigma}-1}{1-\sigma} + \mu\bigl[r a + w - c\bigr].
\]

\textbf{Step 3: First-order conditions.}

\begin{align*}
    \frac{\partial\mathcal{H}}{\partial c} = 0: &\quad c^{-\sigma} = \mu. \\
    \dot{\mu} = \rho\mu - \frac{\partial\mathcal{H}}{\partial a}: &\quad
    \dot{\mu} = \rho\mu - \mu r.
\end{align*}

Differentiating $\mu = c^{-\sigma}$ with respect to time gives
$\dot{\mu} = -\sigma c^{-\sigma-1}\dot{c}$.  Substituting into the costate equation:
\[
    -\sigma c^{-\sigma-1}\dot{c} = (\rho - r)\,c^{-\sigma}
    \implies
    \boxed{\frac{\dot{c}}{c} = \frac{r - \rho}{\sigma}.}
\]

\textbf{Step 4: Why $\tau$ does not appear.}

The tax $\tau$ is absent because the lump-sum rebate $T_t = \tau w_t$ exactly offsets the
labour income tax: the household's disposable income is unchanged. This is
\emph{Ricardian equivalence for lump-sum rebates}: any revenue the government takes from
one hand, it returns through the other, leaving the intertemporal consumption decision
unaltered. Only the after-tax return to capital $r_t$ governs the Euler equation, and
that is untouched by the tax.
\end{answerbox}

\subsection*{2.\ (7 points)}

\begin{answerbox}
\textbf{Steady-state conditions.}

On the BGP, $\dot{c} = 0$ and $\dot{k} = 0$.

\textit{Condition 1 (Modified Golden Rule).} Setting $\dot{c}/c = 0$ in the Euler equation:
\[
    r^* = \rho \implies f'(k^*) - \delta = \rho \implies
    \boxed{f'(k^*) = \rho + \delta.}
\]

\textit{Condition 2 (Market clearing / resource constraint).} With $n = 0$ and no public
spending, capital market clearing requires:
\[
    \dot{k} = f(k) - \delta k - c = 0 \implies \boxed{c^* = f(k^*) - \delta k^*.}
\]

\textbf{Comparison to no-tax steady state.}

The modified Golden Rule $f'(k^*) = \rho + \delta$ is \emph{identical} to the no-tax
Ramsey condition, since $\tau$ does not enter either equation. Therefore $k^*$ is
\emph{unchanged} by the labour income tax with lump-sum rebate.

\textbf{Comparison to the Golden Rule.}

The Golden Rule satisfies $f'(k^{GR}) = \delta$ (with $n=0$). Since $\rho > 0$, we have
$f'(k^*) = \rho + \delta > \delta = f'(k^{GR})$, and by diminishing returns
$k^* < k^{GR}$. The Ramsey economy is always below the Golden Rule (it is impatient).

\textbf{Is the labour income tax distortionary?}

No. Because the tax and the lump-sum rebate cancel in the budget constraint, the tax does
not alter the intertemporal price of consumption (governed only by $r$), nor the
steady-state capital stock. The allocation of resources is identical to the no-tax economy.
The labour income tax with full lump-sum rebate is \emph{non-distortionary}.
\end{answerbox}

\subsection*{3.\ (6 points)}

\begin{answerbox}
\textbf{Transition dynamics.}

Since the economy was already at the no-tax Ramsey steady state $(k^*, c^*)$, and since
the introduction of $\tau > 0$ with a full lump-sum rebate leaves \emph{both} the Euler
equation and the resource constraint unchanged (as shown in parts 1 and 2), the steady
state $(k^*, c^*)$ is unaffected.

On impact at $t=0$: \emph{nothing changes}. Consumption does not jump; capital does not
change (it is a state variable). The household's budget set is identical to what it was
before.

Along the transition path: there \emph{is} no transition. The economy was at the
(unchanged) steady state and remains there.

\textbf{Phase diagram interpretation.}

The phase diagram has two loci:
\begin{itemize}[itemsep=2pt]
    \item $\dot{c}=0$ locus: defined by $f'(k) = \rho + \delta$. Vertical line at $k^*$.
    \item $\dot{k}=0$ locus: defined by $c = f(k) - \delta k$. Hump-shaped.
\end{itemize}
Both loci are determined solely by $f$, $\delta$, and $\rho$ --- none of which depend on
$\tau$. The tax leaves the phase diagram completely unchanged. The economy sits at the
intersection point and stays there.
\end{answerbox}

\subsection*{4.\ (6 points)}

\begin{answerbox}
\textbf{Change to budget constraint.}

If revenue finances a public good $G_t = \tau w_t$ with no utility or production value,
the household no longer receives the lump-sum transfer. The budget constraint becomes:
\[
    \dot{a}_t = r_t a_t + (1-\tau)w_t - c_t.
\]
The term $\tau w_t$ is now absent from disposable income.

\textbf{Effect on the Euler equation.}

The current-value Hamiltonian is now:
\[
    \mathcal{H} = \frac{c^{1-\sigma}-1}{1-\sigma} + \mu\bigl[r a + (1-\tau)w - c\bigr].
\]
The FOC with respect to $c$ still gives $c^{-\sigma} = \mu$, and the costate equation
still gives $\dot{\mu} = (\rho - r)\mu$. Therefore:
\[
    \frac{\dot{c}}{c} = \frac{r - \rho}{\sigma}.
\]
The Euler equation is \emph{unchanged}. The $\dot{c} = 0$ locus is still $f'(k^*) = \rho + \delta$,
so the steady-state capital stock $k^*$ is still the same.

\textbf{Effect on steady-state consumption.}

The resource constraint (economy-wide) is:
\[
    \dot{k} = f(k) - \delta k - c - G = f(k) - \delta k - c - \tau w.
\]
In steady state $\dot{k} = 0$:
\[
    c^* = f(k^*) - \delta k^* - \tau w^*,
\]
where $w^* = f(k^*) - k^* f'(k^*)$. Consumption is \emph{lower} by $\tau w^*$ compared
to the no-tax case, since government spending crowds out private consumption
one-for-one.

\textbf{Is the tax distortionary?}

In terms of capital accumulation: no. The $k^*$ is unchanged because the Euler equation
and the $\dot{c}=0$ locus are the same. The tax does not distort the intertemporal
margin between present and future consumption (which is governed only by $r$).

However, the tax \emph{does} reduce household welfare: steady-state consumption falls by
$\tau w^* > 0$, and on the new BGP households are worse off. The tax distorts the
consumption level (a level effect) but not the capital allocation or growth rate. In this
narrow sense it is a \emph{non-distortionary} tax on allocative efficiency, but it
imposes a welfare cost through the reduction in private consumption.

\textbf{Transition dynamics.}

On impact at $t=0$, the tax is implemented. The steady state shifts from
$(k^*, f(k^*)-\delta k^*)$ to $(k^*, f(k^*)-\delta k^* - \tau w^*)$: the $\dot{k}=0$
locus shifts down by $\tau w^*$ at every $k$. The $\dot{c}=0$ locus is unchanged.
On impact, $c$ must jump \emph{down} to the new saddle path leading to the lower
consumption steady state. Capital $k$ remains at $k^*$ on impact (state variable), then
remains there (since $k^*$ is unchanged). So the transition is instantaneous: $c$ jumps
down to the new $c^*$ immediately, and no further adjustment is needed.
\end{answerbox}

\clearpage

% ═════════════════════════════════════════════════════════════════════════════
\section{Question 2 --- Aghion--Howitt Quality Ladders}
% ═════════════════════════════════════════════════════════════════════════════

\subsection*{1.\ (5 points)}

\begin{answerbox}
\textbf{Step 1: Final-good firm's demand for intermediates.}

The final-good firm maximises profits taking $q_t$ and $p$ as given:
\[
    \max_x \; A q_t x^{\alpha} - p\,x.
\]
First-order condition:
\[
    \alpha A q_t x^{\alpha-1} = p \implies x = \left(\frac{\alpha A q_t}{p}\right)^{1/(1-\alpha)}.
\]

\textbf{Step 2: Monopoly price.}

The current quality leader produces the intermediate good one-for-one from labour at
marginal cost $w$ (wage). With limit pricing against the previous-generation technology
(which would set price equal to $w/\alpha$ to deter entry by the previous leader at quality
$q_t/\lambda$), the unconstrained monopoly price is:
\[
    p^* = \frac{w}{\alpha},
\]
which is a markup $1/\alpha > 1$ over marginal cost $w$.

\textbf{Step 3: Monopoly profits.}

Profits are revenue minus cost:
\[
    \pi = (p^* - w)\,x^* = \left(\frac{w}{\alpha} - w\right)x^*
    = \frac{1-\alpha}{\alpha}\,w\,x^*.
\]
Substituting $x^* = (\alpha^2 A q_t/w)^{1/(1-\alpha)}$:
\[
    \boxed{\pi = \frac{1-\alpha}{\alpha}\,w\,\left(\frac{\alpha^2 A q_t}{w}\right)^{1/(1-\alpha)}.}
\]
Note that $\pi$ is proportional to $q_t$, consistent with BGP scaling.
\end{answerbox}

\subsection*{2.\ (7 points)}

\begin{answerbox}
\textbf{Step 1: Bellman equation for the quality leader.}

Being the current quality leader generates flow profits $\pi$ but the leader is displaced
at Poisson rate $\iota$ (the aggregate rate of creative destruction). The no-arbitrage
Bellman equation is:
\[
    r V = \pi - \iota V + \dot{V},
\]
or equivalently $(r + \iota)V = \pi + \dot{V}$.

\textbf{Step 2: BGP value.}

On the BGP, $r$, $\iota$, and $\pi$ are constant. Conjecture $\dot{V} = 0$ (the value of
leadership is constant in BGP units). Then:
\[
    \boxed{V^* = \frac{\pi}{r + \iota}.}
\]
The value equals the discounted flow of profits, with $r + \iota$ as the effective
discount rate (combining time preference and the risk of displacement).

\textbf{Step 3: Free-entry condition.}

Each unit of research labour generates a successful innovation at rate $\lambda$. A
successful innovation delivers value $V^*$ (the new leader captures the prize). The
expected return per unit of research labour per unit of time is $\lambda V^*$. The cost
of one unit of research labour is the wage $w$. Free entry drives expected profits to zero:
\[
    \lambda V^* = w.
\]

\textbf{Step 4: Equilibrium creative destruction rate.}

Substituting $V^* = \pi/(r+\iota)$ into the free-entry condition:
\[
    \frac{\lambda \pi}{r + \iota} = w
    \implies r + \iota = \frac{\lambda\pi}{w}
    \implies \iota^* = \frac{\lambda\pi}{w} - r.
\]
Using $\pi = \frac{1-\alpha}{\alpha}w x^*$:
\[
    \boxed{\iota^* = \lambda\,\frac{1-\alpha}{\alpha}\,x^* - r,}
\]
where $x^* = L - L_R = L - \iota^*/\lambda$ gives a fixed-point equation in $\iota^*$.
This pins down the equilibrium jointly with the labour market clearing condition.
\end{answerbox}

\subsection*{3.\ (5 points)}

\begin{answerbox}
\textbf{Business-stealing externality.}

When an entrant researcher successfully innovates, she displaces the incumbent quality
leader. From the incumbent's perspective, this is a loss of $V^*$ (the value of
leadership vanishes). From the social planner's perspective, this transfer of leadership
is a pure redistribution of monopoly rents --- no new real resources are created by
displacement itself. The entrant ignores the loss she imposes on the incumbent and
focuses only on the prize $V^*$ she gains. This leads researchers to over-invest in R\&D
relative to what a social planner would choose. \textbf{Direction: pushes $\iota$ above
the social optimum.}

\textbf{Appropriability (consumer-surplus) externality.}

A quality improvement raises the productivity of the intermediate good, generating
surplus for all downstream final-good producers and ultimately consumers. The
innovating monopolist, however, captures only her monopoly profits $\pi$, which are a
fraction of the total social value of the innovation. The rest --- the consumer surplus
triangle between the competitive price $w$ and the monopoly price $w/\alpha$ ---
goes uncaptured by the innovator. Because the innovator cannot appropriate the full
social return, private incentives to innovate are too low. \textbf{Direction: pushes
$\iota$ below the social optimum.}
\end{answerbox}

\subsection*{4.\ (4 points)}

\begin{answerbox}
\textbf{Condition for $\iota^* > \iota^{\text{opt}}$.}

Equilibrium R\&D exceeds the social optimum when the \textbf{business-stealing
externality dominates} the appropriability externality. This occurs when the value
destroyed by displacing the incumbent ($V^*$, lost to the old leader) is large relative
to the consumer surplus that the innovator fails to capture.

Formally, define the social value of innovation as $W^* = V^* + CS^*$, where $CS^*$ is
the consumer surplus gain from the quality improvement. The social planner equates the
marginal cost of research to the full social return $\lambda W^*$. The decentralised
equilibrium equates it to $\lambda V^*$. Over-investment occurs when:
\[
    \lambda V^* > \lambda W^* \iff V^* > W^*,
\]
which requires the loss imposed on the incumbent ($V^*$, internalised by the social
planner as a cost) to outweigh the consumer surplus gain. Since $V^* = \pi/(r+\iota)$,
this is more likely when monopoly profits $\pi$ are large relative to the consumer
surplus --- i.e., when the markup $1/\alpha$ is large (small $\alpha$) and the step size
$\lambda$ is not too large (so the consumer surplus gain is modest).
\end{answerbox}

\subsection*{5.\ (4 points)}

\begin{answerbox}
\textbf{Policy mix to implement the social optimum.}

Two externalities operate in opposite directions, so a single instrument cannot correct
both simultaneously. A proposed mix:

\textbf{Instrument 1: R\&D subsidy $s_R$.} A per-unit subsidy to research labour raises
the private return to innovation, correcting the appropriability externality (which makes
private $\iota$ too low). The subsidy effectively raises the private prize from $\lambda
V^*$ to $\lambda V^* + s_R$, encouraging more R\&D.

\textbf{Instrument 2: Innovation tax / patent narrowing.} A tax on the profits of
successful innovators (or equivalently, patent policies that shorten the effective
monopoly period) reduces $V^*$, correcting the business-stealing externality (which makes
private $\iota$ too high). By reducing the prize for displacement, this instrument
internalises the rent-destruction cost.

\textbf{Why a single instrument fails.} A pure R\&D subsidy would correct appropriability
but exacerbate business-stealing. A pure R\&D tax would reduce business-stealing but
worsen the appropriability problem. The social optimum requires calibrating both: the
subsidy to close the appropriability gap and the tax (or patent design) to close the
business-stealing gap simultaneously. Patent breadth and duration policies that
independently affect $\pi$ and $V^*$ are a natural two-dimensional instrument set.
\end{answerbox}

\clearpage

% ═════════════════════════════════════════════════════════════════════════════
\section{Question 3 --- Development Accounting with Heterogeneous Human Capital}
% ═════════════════════════════════════════════════════════════════════════════

\subsection*{1.\ (6 points)}

\begin{answerbox}
\textbf{Step 1: Divide by $L_i$.}

\[
    \frac{Y_i}{L_i} = \frac{K_i^{\alpha}(h_i L_i)^{1-\alpha} A_i}{L_i}
    = K_i^{\alpha}\,h_i^{1-\alpha}\,L_i^{1-\alpha-1}\,A_i
    = K_i^{\alpha}\,L_i^{-\alpha}\,h_i^{1-\alpha}\,A_i.
\]

\textbf{Step 2: Re-express capital in terms of the capital-output ratio.}

Write $K_i = (K_i/Y_i)\cdot Y_i$. Then:
\[
    K_i^{\alpha} = \left(\frac{K_i}{Y_i}\right)^{\!\alpha} Y_i^{\alpha}.
\]
Substituting and using $y_i \equiv Y_i/L_i$:
\[
    y_i = \left(\frac{K_i}{Y_i}\right)^{\!\alpha} y_i^{\alpha}\,h_i^{1-\alpha}\,A_i.
\]
Collect $y_i^{\alpha}$ on the right and $y_i$ on the left:
\[
    y_i^{1-\alpha} = \left(\frac{K_i}{Y_i}\right)^{\!\alpha} h_i^{1-\alpha} A_i
    \implies
    y_i = \left(\frac{K_i}{Y_i}\right)^{\!\alpha/(1-\alpha)} h_i\,A_i.
\]

\textbf{Development accounting identity:}
\[
    \boxed{y_i = \underbrace{\left(\frac{K_i}{Y_i}\right)^{\!\alpha/(1-\alpha)}}_{\text{capital intensity}}
    \cdot\underbrace{h_i}_{\text{human capital}}
    \cdot\underbrace{A_i}_{\text{TFP}}.}
\]
With $\alpha = 1/3$, the exponent $\alpha/(1-\alpha) = 1/2$.
\end{answerbox}

\subsection*{2.\ (8 points)}

\begin{answerbox}
\textbf{Step 1: Total log income gap.}

\[
    \ln\!\left(\frac{y_{US}}{y_E}\right) = \ln\!\left(\frac{1.00}{0.10}\right) = \ln 10 \approx 2.303.
\]

\textbf{Step 2: Contribution of physical capital intensity.}

With $\alpha/(1-\alpha) = 1/2$:
\[
    \Delta\ln y_{\text{capital}} = \frac{1}{2}\ln\!\left(\frac{\kappa_{US}}{\kappa_E}\right)
    = \frac{1}{2}\ln\!\left(\frac{3.00}{2.00}\right)
    = \frac{1}{2}\ln(1.5) = \frac{1}{2}(0.405) \approx 0.203.
\]
Share of gap: $0.203 / 2.303 \approx \mathbf{8.8\%}$.

\textbf{Step 3: Contribution of human capital.}

\[
    \Delta\ln y_{\text{human}} = \ln\!\left(\frac{h_{US}}{h_E}\right)
    = \ln\!\left(\frac{1.00}{0.55}\right) = \ln(1.818) \approx 0.598.
\]
Share of gap: $0.598 / 2.303 \approx \mathbf{26.0\%}$.

\textbf{Step 4: TFP residual.}

\[
    \Delta\ln y_{\text{TFP}} = 2.303 - 0.203 - 0.598 = 1.502.
\]
Share of gap: $1.502 / 2.303 \approx \mathbf{65.2\%}$.

\textbf{Summary table.}

\begin{center}
\begin{tabular}{lcc}
\toprule
Source & Log contribution & Share of gap \\
\midrule
Capital intensity $\kappa^{1/2}$ & 0.203 & 8.8\% \\
Human capital $h$                & 0.598 & 26.0\% \\
TFP residual $A$                 & 1.502 & 65.2\% \\
\midrule
Total                            & 2.303 & 100\% \\
\bottomrule
\end{tabular}
\end{center}

TFP accounts for roughly \textbf{65\%} of the log income gap, consistent with the
Hall--Jones (1999) finding that TFP differences dominate cross-country income variation.
\end{answerbox}

\subsection*{3.\ (6 points)}

\begin{answerbox}
\textbf{The misallocation critique.}

The development accounting exercise attributes the TFP residual $A_i$ to ``technology.''
Hsieh and Klenow (2009) argue that a substantial portion of this residual instead reflects
\emph{resource misallocation}: capital and labour are not allocated to their most
productive uses across firms within an economy.

\textbf{TFPR and its equalization under no distortions.}

In the Hsieh--Klenow framework, each firm $s$ has physical TFP (TFPQ$_s$) but also faces
firm-specific output wedges $\tau_{Y,s}$ and capital wedges $\tau_{K,s}$. Revenue TFP is:
\[
    \mathrm{TFPR}_s \equiv P_s \cdot \mathrm{TFPQ}_s.
\]
Under perfect competition with no distortions, profit maximisation with a common markup
equates TFPR across all firms: higher-productivity firms charge lower prices so that
$P_s \cdot \mathrm{TFPQ}_s$ is the same for all $s$. Dispersion in TFPR is therefore a
signal that distortions are present.

\textbf{How TFPR dispersion reduces aggregate TFP.}

When wedges prevent resources from flowing to high-TFPQ firms, aggregate TFP falls below
its frictionless benchmark. Hsieh and Klenow show:
\[
    \frac{A_i}{A_i^*} = \prod_s \left(\frac{\overline{\mathrm{TFPR}}}{\mathrm{TFPR}_s}\right)^{\text{share}_s} < 1
    \quad \text{whenever TFPR is dispersed.}
\]
High-productivity firms are constrained to operate below efficient scale; low-productivity
firms survive and absorb resources they would not in a frictionless equilibrium. The
cross-sectional standard deviation of log TFPR quantifies the severity of this
misallocation.

\textbf{Implication for development accounting.}

The TFP residual $A_i$ in the standard decomposition conflates true technology differences
with misallocation losses. Two countries could have identical production technologies at
the firm level yet exhibit large measured TFP gaps if one suffers from severe misallocation.
Policies that reduce wedges (entry barriers, credit distortions, product market regulations)
can raise measured aggregate TFP without any technology transfer.
\end{answerbox}

\subsection*{4.\ (5 points)}

\begin{answerbox}
\textbf{Implication: policy should target the sources of wedges, not just accumulation.}

Since TFP (including misallocation) explains most of the income gap, and since much of
the TFP gap may reflect misallocation rather than technology, development policy should
focus on reducing the distortions that cause TFPR dispersion --- not just subsidising
investment in physical or human capital.

\textbf{Policy recommendation 1: Reduce barriers to firm entry, exit, and reallocation.}

Entry barriers (licensing, regulatory costs) prevent high-productivity entrants from
displacing low-productivity incumbents. Exit barriers (bail-outs, labour regulations that
prevent layoffs) keep low-productivity firms alive. Removing these barriers allows
resources to flow toward high-TFPQ firms, raising aggregate TFP through reallocation.
\textit{Limitation}: identifying the binding barrier requires firm-level data. Reforms
that open entry may also increase exit, causing short-run disruption and political
resistance. The appropriate speed of liberalisation is context-dependent.

\textbf{Policy recommendation 2: Deepen and reform financial markets.}

Capital wedges $\tau_{K,s}$ often arise because high-productivity firms --- frequently
small, young, or informal --- cannot borrow at the same terms as large established firms.
Strengthening creditor rights, improving information asymmetry (credit registries), and
developing equity markets enables productive firms to access capital and expand.
\textit{Limitation}: credit market deepening can also fuel misallocation if lenders
themselves have poor information (financing connected firms over productive ones).
Financial sector reforms require complementary institutional improvements (rule of law,
accounting standards) to be effective, making them slow and difficult to implement.
\end{answerbox}

\clearpage

% ═════════════════════════════════════════════════════════════════════════════
\section{Short Questions}
% ═════════════════════════════════════════════════════════════════════════════

\subsection*{SQ1.\ (8 points) Kaldor Facts and Balanced Growth}

\begin{answerbox}
\textbf{Three Kaldor stylised facts.}

\begin{enumerate}[label=(\roman*), itemsep=3pt]
    \item The growth rate of output per worker is roughly constant over long periods.
    \item The capital-output ratio $K/Y$ is roughly constant over time.
    \item The return to capital $r$ is roughly constant over time, and labour and capital
    shares of income are roughly constant.
\end{enumerate}
(Additional Kaldor facts include persistent cross-country differences in growth rates and
the positive correlation between the investment rate and output per worker.)

\textbf{Which production functions are consistent with balanced growth?}

Balanced growth requires constant growth rates of all per-worker quantities and constant
factor shares simultaneously. Uzawa (1961) showed that balanced growth in a
neoclassical model requires that technical progress be \emph{labour-augmenting}
(Harrod-neutral). Given Harrod-neutral technology, the production function must satisfy
certain conditions on the elasticity of substitution. Specifically, constant factor
shares on a BGP require that the elasticity of substitution between capital and
labour-augmenting technology converges to 1 as the economy grows --- which is exactly the
property of \emph{Cobb-Douglas} production functions, and more generally of functions
that approach Cobb-Douglas asympotically (such as CES with unit elasticity in the limit).

\textbf{The privileged role of Cobb-Douglas.}

The Cobb-Douglas function $Y = K^{\alpha}(AL)^{1-\alpha}$ delivers constant factor
shares \emph{exactly and for all input ratios}, not merely asymptotically. Capital's
share is always $\alpha$ regardless of the capital-labour ratio, and the labour share is
always $1-\alpha$. This makes Cobb-Douglas the unique production function that is
simultaneously consistent with all Kaldor facts at every point along the transition path,
not only on the BGP. It is therefore the workhorse of both growth theory and empirical
work: it is simple, tractable, and exactly consistent with the long-run regularities
identified by Kaldor.
\end{answerbox}

\subsection*{SQ2.\ (9 points) TVC and the No-Ponzi Condition}

\begin{answerbox}
\textbf{The Transversality Condition (TVC).}

The TVC for the household's optimal control problem in the Ramsey model (with costate
variable $\mu_t$ = shadow price of wealth, discount $\rho$, capital $k_t$) is:
\[
    \lim_{t\to\infty} e^{-\rho t}\,\mu_t\,k_t = 0.
\]
Equivalently, since $\mu_t = u'(c_t)$ and using the transversality interpretation in
terms of asset value:
\[
    \lim_{t\to\infty} e^{-\int_0^t r_s\,\dd s}\,k_t = 0.
\]
Economically: the present value of the terminal capital stock must be zero. The household
should not arrive at ``infinity'' with assets that still have positive value --- it would
have been optimal to consume more instead.

\textbf{The No-Ponzi Condition (NPC).}

The NPC constrains the household to not run explosive debt:
\[
    \lim_{t\to\infty} a_t\,e^{-\int_0^t r_s\,\dd s} \geq 0.
\]
The household cannot borrow more than the present value of its future repayments. Without
this constraint, households could finance unlimited current consumption by rolling over
debt forever (a Ponzi scheme). The NPC closes the household's optimisation problem.

\textbf{How the TVC rules out over-accumulation.}

Suppose the economy follows a path with $k_t \to \infty$. As $k$ grows without bound,
$f'(k) \to 0 < \rho + \delta$, so the capital return falls below the discount rate.
Along such a path, the costate $\mu_t = u'(c_t)$ must be growing fast enough to satisfy
the costate equation $\dot{\mu} = (\rho - r)\mu > 0$. This means $\mu_t$ grows
explosively and $e^{-\rho t}\mu_t k_t \to \infty$, violating the TVC. The TVC therefore
rules out all over-accumulation paths as sub-optimal: the household would do better to
consume more and let capital fall back toward $k^*$.

\textbf{Contrast with the Solow model.}

In the Solow model, the saving rate $s$ is fixed exogenously. If $s$ is set above the
Golden Rule saving rate $s^{GR} = \alpha$ (Cobb-Douglas), the economy accumulates capital
beyond $k^{GR}$ and enters the dynamically inefficient region where $f'(k) < \delta + n$.
There is no optimising household to impose the TVC, and no mechanism to prevent
over-accumulation. The Solow model therefore admits dynamically inefficient equilibria as
possible steady states, while the Ramsey model does not.
\end{answerbox}

\subsection*{SQ3.\ (8 points) Baumol Cost Disease and Aggregate Productivity}

\begin{answerbox}
\textbf{What is Baumol's cost disease?}

Baumol's cost disease (Baumol and Bowen, 1966) refers to the phenomenon by which the
relative price of labour-intensive services rises persistently over time, and the
employment share of these services tends to increase, even without any change in their
productivity.

\textbf{Mechanism: why the stagnant sector's relative price rises.}

In a two-sector economy with a progressive sector (productivity grows at rate $g > 0$)
and a stagnant sector (productivity fixed), a competitive labour market equalises wages
across sectors at $w_t = A_{P,t}$. Since the stagnant sector's productivity does not
grow, the real cost of producing a unit of stagnant output rises at rate $g$: a unit of
stagnant output requires the same amount of labour as always, but that labour now costs
$w_t = A_0 e^{gt}$. The relative price $p_S/p_P$ rises at rate $g$, reflecting the rising
opportunity cost of the labour that must be devoted to the stagnant sector.

\textbf{What drives rising employment in the stagnant sector.}

For a given expenditure share on stagnant services, constant productivity in that sector
means that more labour is required to produce the demanded quantity as wages rise. If
preferences are non-homothetic with an income elasticity of stagnant goods exceeding one,
rising income further increases the demand for services, amplifying the employment shift.

\textbf{Does rising service employment imply lower aggregate productivity growth?}

Not necessarily. In the Baumol two-sector model, aggregate productivity growth equals the
weighted average of sectoral productivity growth rates. If the employment share of the
stagnant sector rises, this weight shifts toward zero and aggregate productivity growth
falls over time. However, under \textbf{Cobb-Douglas preferences}, expenditure shares are
constant, which implies employment shares are also constant (since both sectors' outputs
are produced with labour alone). In this case, aggregate productivity growth equals
$\beta g$ (the expenditure-weighted average), and it remains constant over time regardless
of the rising relative price of services. Rising service \emph{employment} as a share of
total employment does not materialise under Cobb-Douglas preferences precisely because
the unit-elastic substitution keeps spending shares pinned, and the labour allocation
tracks spending.
\end{answerbox}

\vspace{20pt}
\noindent\rule{\textwidth}{0.4pt}

\begin{center}
    \textit{End of Practice Exam 5 Answer Key.}
\end{center}

\end{document}
